(1) 自由传播$\mathsf P(L)=\begin{pmatrix}1&L\\0&1\end{pmatrix}$，凹面镜的近轴反射等效$\mathsf M(R)=\begin{pmatrix}1&0\\-2/R&1\end{pmatrix}$。从平面镜起算往返矩阵$\mathsf A=\mathsf P(L)\mathsf M(R)\mathsf P(L)$，其迹$\operatorname{tr}\mathsf A=2(1-2L/R)$、行列式为1。两个光线本征值在单位圆的严格条件为$|\operatorname{tr}\mathsf A|<2$，即$0<L<R$；在$L=0,R$为边界中性而非严格稳定。
(2) 平面镜处高斯模波前曲率无穷，故它恰位于束腰，取$q(0)=iz_R$。传播到凹镜后波前几何曲率$R_{\rm wave}(L)=L[1+(z_R/L)^2]$必须与镜面曲率$R$相等，故$z_R^2=L(R-L)>0$。由$z_R=\pi w_0^2/\lambda$得到$w_0^2=(\lambda/\pi)\sqrt{L(R-L)}$；在凹镜处束宽$w(L)=w_0\sqrt{1+(L/z_R)^2}$，与谐振腔口径匹配。
(3) Gouy相位由束腰到凹镜累积$\psi=\arctan(L/z_R)=\arccos\sqrt{1-L/R}$。一次往返的纵向相位为$2kL$，横向阶数$N=m+n$的Gouy相位为$2(N+1)\psi$；回程共振条件$2kL-2(N+1)\psi=2\pi q$，故$\nu_{qmn}=(c/2L)[q+(m+n+1)\psi/\pi]$。横向同阶$m+n$的不同模式理想简并，实际镜像散、像差使简并裂开。
(4) 若$L=R/2$，则$z_R=R/2$且$\psi=\pi/4$，部分横模与不同纵模可能近简并；当$L\to0$时$z_R\sim\sqrt{LR}$，严格近轴高斯模的物理孔径和衍射损耗必须另核查；当$L\to R^-$，$z_R\to0$，凹镜面光束发散，进入平面—凹镜构型的稳定性边缘。此处不是两凹镜“同心腔”的几何条件，不能混同。
